\scanpage{40}
Considérons l'espèce de structure ``anneau (comm. unifère)'', soit $T$.
On sait qu'elle est définie en termes de produits finis
soit $\lambda_0$ \inserted{$=(\text{ensembles finis},\varnothing)$} le type de structure correspondant. \deleted{On se propose de \unsure{$\xi$}, étudier}
\inserted{Donc $T$ définit un $\lambda_0$-type $T_{\lambda_0}$,}
donc pour tout $\lambda$ type tel que $\lambda_0\longrightarrow\lambda$ (les
$\lambda$-catégories ont des produits finis, les $\lambda$-foncteurs y commutent)
il définit un $\lambda$-type $T_\lambda$. On se propose, pour divers
$\lambda$, d'étudier l'existence \inserted{et l'unicité} d'un sous-type plein $T_\lambda^*$
($T_\lambda^*(C)\mathbin{\unsure{-|\to}}T_\lambda(C)$ pl. fidèle pour tout $C$) de telle façon que
pour $C=(\mathrm{Ens})$, $T_\lambda^*(C)$ soit la sous-catégorie \inserted{str.} pleine de
$T_\lambda(C)$ \inserted{$=(\mathrm{Ann})$} $=$ cat. des anneaux (comm. unifères) qui sont des \uline{corps} \unsure{$\xi$}
(commutatifs). [On sait que l'unicité est acquise si $\lambda=(\underset{\leftarrow}{\lambda}_1,\underset{\longrightarrow}{\lambda}_2=\varnothing)$].

a) Si $\lambda=\lambda_0$, \uline{il n'existe pas} une telle $T_\lambda$. En effet, \deleted{\unsure{en termes}} \deleted{\unsure{$\lambda$ produit}}
si $\lambda=(\underset{\leftarrow}{\lambda}_1,\underset{\rightarrow}{\lambda}_2=\varnothing)$, $\underset{\leftarrow}{\lambda}$ formé de \unsure{diagrammes|diag. non} finis, alors
$T_\lambda$ est $\lambda$-\unsure{représenté|représentable} par la $\lambda$-catégorie \inserted{$\mathcal R_{\lambda}^{T}$} qui est \inserted{l'opposé de} la sous-catégorie pleine $S_\lambda^T$ de Ann engendrée (pour les $\varinjlim$ de type $\unsure{\lambda_1|\lambda_3}$)
par $\mathbb Z[T]$ (l'objet libre engendré par 1 élément). Les sous-\deleted{catégories}types pleins de $T_\lambda$ \deleted{sont les} correspondent aux
catégories de fractions de $\mathcal R_\lambda^T$ par des ens. de flèches $M^*$
stables par composition, et par $\varprojlim$ de type $\lambda_1$ (ou aux
cat. \deleted{\unsure{quot.|quotient}} de fractions à \uline{gauche} pour $S_\lambda^T$ par des ens
$M$ de flèches stables par $\varinjlim$ de type $\lambda_1$). Ici, \deleted{des}\inserted{l'ens $M$}
flèches en question \deleted{\unsure{sont}}\inserted{est} \deleted{celles}\inserted{formé} des flèches $u:A\mathbin{\unsure{-|\to}}B$
dans $S_\lambda^T$ \inserted{$\subset(\mathrm{Ann})$} telles que pour tout \deleted{\unsure{k}} \uline{corps} $k\in\mathrm{Ob}(\mathrm{Ann})$,
la flèche $\mathrm{Hom}_{\mathrm{ann}}(B,k)\longrightarrow\mathrm{Hom}_{\mathrm{ann}}(A,k)$ soit \uline{bijective}, [ou
encore que le \deleted{\unsure{hom}} morphisme de schémas affines $u':X\mathbin{\unsure{-|\leftarrow|\to}}\deleted{Y}$
défini par $u$ soit tel que $X(k)\xrightarrow{\sim}Y(k)$ pour tt corps $k$.
On sait que cela signifie que $u$ est \uline{bijective} et définit
des ext. résiduelles triviales] \deleted{\unsure{Alors}} Cet ens de flèches
est évidemment stable par composition et par $\varinjlim$ (resp. $\varprojlim$)
de type $\lambda_1$ (serait vrai pour tte \unsure{sous|ss}-catégorie de (Ann),
% The ruled sheet is upside down in the scan; long oblique pen stroke through the working.
% archive sheet number: 39
